%% kapa: exam=1 round=37 year=2026 period=1 form=A subject=economics subject_ko=경제학원론 seq=14 number=54
%% source: Q-Net 감정평가사 기출문제 (cst00309) artlSeq=5258924 file=raw/1008/5258924/2026년 제37회 감정평가사 1차 1교시.pdf
%% conversion: kordoc 4.24.51 -> tools/gen1.py (auto); review=proofread-sonnet; flags=none
\begin{question}{54}
파레토 효율성과 관련된 설명 중 옳지 않은 것은? (단, $X$, $Y$ 두 재화 / A, B 2인 / $L$(노동), $K$(자본) 두 생산요소 경제를 가정한다.)
\begin{choices}
  \item 소비의 효율성은 두 재화에 대한 $A$와 $B$의 한계대체율($MRS_{XY}$)이 같을 때 달성된다.
  \item 생산의 효율성은 두 생산요소에 대한 $X$와 $Y$의 한계기술대체율($MRTS_{LK}$)이 같을 때 달성된다.
  \item 생산의 효율성 조건을 만족하는 점들을 $L-K$평면으로 옮겨 놓은 것을 생산가능곡선이라 한다.
  \item 생산의 효율성 조건을 만족하는 점들을 연결한 선을 생산의 계약곡선이라 한다.
  \item 생산의 효율성, 소비의 효율성, 경제 전체의 효율성을 모두 만족하면 파레토 효율성이 달성된다.
\end{choices}
\end{question}
\answer{3}
